% =====================================================================
%  Figure environments.
%
%  Panels 1-10 are composed by Paper-materials/figs.py --panels from the
%  individual figures in Paper-materials/figures/.  Figures 11-14 come from
%  the network pipeline (PYTHON_Modelling_Fefet_Codes/ecg detection/fig_k*.py)
%  and are copied into Paper-materials/figures/ unchanged.
%
%  Every panel's constituent figures, their source runs, the claim each one
%  supports and the normalization each used are recorded in
%  Paper-materials/figures_manifest.csv.
% =====================================================================

\begin{figure*}[t]
\centering\includegraphics[width=\textwidth]{panel01.png}
\caption{Structure and normalization. (a) The simulated cross-section to scale: a
\SI{5}{\nano\metre} silicon nanosheet body gated on both faces through
\SI{1}{\nano\metre} \ce{SiO2} and \SI{7}{\nano\metre} HZO, with a
\SI{100}{\nano\metre} gate and \SI{15}{\nano\metre} source and drain overlaps. The
marked point is where polarization and electric field are probed in every transient
reported here. (b) The mapping from the two-dimensional slab to the nanosheet. Because
the slab is gated on both faces, a slab of depth $A$ carries gate width $2A$; matching the
nanosheet perimeter $2(W+T_\mathrm{si})=\SI{90}{\nano\metre}$ fixes
$A=\SI{0.045}{\micro\metre}$. Runs executed at \num{0.071}; the correction is an exact
post-multiplier, verified against a re-simulated node to \num{1.8e-16}.}
\label{fig:panel1}
\end{figure*}

\begin{figure*}[t]
\centering\includegraphics[width=\textwidth]{panel02.png}
\caption{Calibration against published measurements of a comparable gate-all-around
FeFET. (a) Overlay in the operating region, with the memory window and both threshold
voltages computed from the plotted curves at the stated constant-current criterion. A
single free width scalar, fitted on the erased branch's saturation, is applied to both
simulated branches; the vertical axis is labelled accordingly. (b) The same data with the
window and the current level divided out --- each curve against its own gate overdrive,
normalized to its own on-current --- so that only the turn-on shape remains. (c) Residual
in decades against the measurement, with a $\pm\num{0.3}$ decade band.}
\label{fig:panel2}
\end{figure*}

\begin{figure*}[t]
\centering\includegraphics[width=\textwidth]{panel03.png}
\caption{The ferroelectric material against what the gate stack reaches. (a) Standalone
metal--ferroelectric--metal capacitor at three sweep amplitudes. Twice the coercive
voltage does not saturate the loop; the saturated amplitude reproduces the calibrated
$P_s$, $|P_r|$ and $F_c$, and the unsaturated amplitudes are retained as the control.
(b) The loop measured through the gate stack, sign-corrected, against the same material's
coercive field. A \SI{\pm4}{\volt} gate sweep reaches about half of $F_c$: this is a deep
minor loop of the stack, not the material loop.}
\label{fig:panel3}
\end{figure*}

\begin{figure*}[t]
\centering\includegraphics[width=\textwidth]{panel04.png}
\caption{DC characteristics. (a) Retained transfer characteristic, both states from the
same run and read after settling. Threshold markers are at the stated constant-current
criterion, the subthreshold slopes are fitted over at least two decades starting a decade
above the branch minimum, and the ambipolar minima are labelled rather than cropped.
(b) Output characteristics for the erased state, the programmed state and intermediate
ladder stops; the inset is the low-drain-bias region, with the computed linear-region
correlation given below the axes.}
\label{fig:panel4}
\end{figure*}

\begin{figure*}[t]
\centering\includegraphics[width=\textwidth]{panel05.png}
\caption{Synaptic behaviour. (a) The potentiation ladder with its measured cycle-to-cycle
band. Filled markers are separable at three standard deviations, open markers are not; the
levels and the noise come from the same eleven-repeat run. (b) The conventional
exponential-saturation nonlinearity fit, with the residual shown as an inset: it is
systematic and one-sided because the train is sigmoidal, so the \numrange{10}{90} percent
pulse span is quoted alongside the fitted parameter. (c) Three potentiation--depression
cycles on a log axis. The window lost between cycles is lost at the floor, not the
ceiling. This is write repeatability, not endurance.}
\label{fig:panel5}
\end{figure*}

\begin{figure*}[t]
\centering\includegraphics[width=\textwidth]{panel06.png}
\caption{Write transients and energy, all from one fifteen-pulse train. (a) Polarization
and gate charge stitched onto one time axis, with the write windows shaded.
(b) Per-pulse polarization transients, colour-graded by pulse index, and the amount
switched per pulse: switching peaks in the middle of the train and falls as the film
saturates. (c) Switching energy per pulse and its cumulative total, taken from the
solver's gate-charge trace rather than from a $V_G I_G$ integral, which at this sampling
interval would measure the sampler.}
\label{fig:panel6}
\end{figure*}

\begin{figure*}[t]
\centering\includegraphics[width=\textwidth]{panel07.png}
\caption{Integrate-and-fire operation. (a) Read current against pulse number at
$V_G=0$, crossing a fixed threshold on the ninth pulse at $18.4\times$ the rest current,
measured against this run's own baseline. (b) Membrane-state accumulation for five drive
amplitudes, measured on the polarization probe during the write; pulse amplitude sets the
integration rate. (c) Gate drive, polarization and read current on one time axis. With the
depolarization time constant used here the device latches after firing, so this is a
single-shot event rather than a repeating leaky cycle.}
\label{fig:panel7}
\end{figure*}

\begin{figure*}[t]
\centering\includegraphics[width=\textwidth]{panel08.png}
\caption{Design space. (a) The coordinate-descent trajectory, six turns, with the value
chosen at each turn. (b) Interfacial oxide thickness: retained window against off-current,
in both evaluation contexts. (c) Ferroelectric thickness as a Pareto choice --- window
against operating voltage, marker area proportional to gate charge. (d) Read bias: the
true retained window and the rest current against $V_\mathrm{read}$, with the leakage
region shaded, and the superseded virgin-referenced dynamic range overlaid as a dashed
line to document the disagreement.}
\label{fig:panel8}
\end{figure*}

\begin{figure*}[t]
\centering\includegraphics[width=\textwidth]{panel09.png}
\caption{Retention and read disturb, all measured after settling. (a) A short hold, with
the depolarization transient and the plateau separated and the wait time stated.
(b) Long retention at two temperatures, for the fully programmed state and for a
mid-ladder analog state, normalized to each plateau. No ten-year extrapolation is drawn
and no activation energy is extracted; a bound over the measured decades is reported
instead. (c) Read disturb against equivalent read count at $V_\mathrm{read}=0$. The
transient is shown in grey so that the distinction between it and disturb is visible
rather than asserted.}
\label{fig:panel9}
\end{figure*}

\begin{figure*}[t]
\centering\includegraphics[width=\textwidth]{panel10.png}
\caption{Comparison and benchmark. (a) The nanosheet against a planar ablation of the
identical stack, on the axes a synapse is judged by; each axis is normalized to the better
of the two devices and the level axis uses one shared separation criterion. The planar
device wins the raw window and loses the application-relevant axes. (b) Analog levels
against energy per programming pulse for fifteen published ferroelectric devices, with
this work marked at both its open-loop and its write-verify level counts. Devices that
report a level count but no energy appear as a marginal rug; open markers are points whose
source number could not be confirmed in the paper text.}
\label{fig:panel10}
\end{figure*}

\begin{figure}[t]
\centering\includegraphics[width=\columnwidth]{K2_accuracy_vs_levels.png}
\caption{Deployed accuracy against the number of conductance levels, for two grid
constructions, with the spread over interior-level placement at fixed count. Accuracy is
not monotonic in level count below eight levels, and the spread at fixed count is larger
than the trend across counts.}
\label{fig:k2}
\end{figure}

\begin{figure}[t]
\centering\includegraphics[width=\columnwidth]{K3a_accuracy_vs_calibration.png}
\caption{Deployed accuracy across twenty corner devices as a function of how much
per-device calibration is allowed: none, one gain scalar, or full per-level write-verify.
The corners are a worst-case envelope rather than a distribution and are never injected as
a single standard deviation.}
\label{fig:k3a}
\end{figure}

\begin{figure}[t]
\centering\includegraphics[width=\columnwidth]{K3c_accuracy_vs_shape_error.png}
\caption{Accuracy after a gain trim against the level misplacement the trim leaves behind.
This is the quantity that tracks accuracy (Spearman \num{-0.85}); the readout-overlap
statistics do not.}
\label{fig:k3c}
\end{figure}

\begin{figure}[t]
\centering\includegraphics[width=\columnwidth]{K3b_accuracy_vs_c2c.png}
\caption{Accuracy against injected cycle-to-cycle noise, for the fifteen-level grid and
for the eight-level grid the three-sigma separability criterion selects. The measured noise
is marked. Noise is injected per conductance branch, not per weight.}
\label{fig:k3b}
\end{figure}

\begin{figure}[t]
\centering\includegraphics[width=\columnwidth]{K4_energy_per_inference.png}
\caption{Energy per heartbeat, with every count taken from the deployed network. The
peripheral conversion bar is drawn on the same axis as the core compute rather than left
to a caveat; the \SI{1}{\pico\joule} per conversion figure is the single assumption in the
budget and is labelled as such.}
\label{fig:k4}
\end{figure}
